% ============================================================
% 读透社·考场前线（杂志页码 02–15）
% 含：热点阅卷观（英国女王专题）+ 10秒鲜素材
% ============================================================
\colchapter{读透社·考场前线}

\section*{英国女王伊丽莎白二世去世：光辉岁月落幕，时代记忆永存}
\auth{热点阅卷观 · 杂志第02–09页}

当地时间2022年9月8日，“超长待机”的英国女王伊丽莎白二世，在苏格兰巴尔莫勒尔堡去世，享年96岁。

伊丽莎白二世于1926年出生，1952年继承王位，在位时间长达70年，历经了15位英国首相，是英国历史上最长寿、在位时间最长的君主。她的人生像是一部厚重的百年史书，经历过第二次世界大战、“冷战”、欧盟诞生与英国入欧脱欧等历史大变革，是那个波谲云诡时代的亲历者和见证者。在位期间，她始终温和、镇定，以持续的个人魅力，化解民众对王室的敌意，成为最受世界欢迎的英国人。

伊丽莎白二世的逝世，是一个时代的落幕——历史，正在随之远去。虽然，我们回顾女王的一生，似乎稍显平淡，说不出她有什么改变人类历史的成就，但也很难想出她做过什么“失礼”的事情，制造出任何引人厌恶的丑闻。而无论未来如何，她都已经完成一位女王的使命，也走完了一个女人履行责任的96年。

\subsection*{人像：英国女王：我需要被人们看到}
\auth{文/达达}

伊丽莎白二世在其70年的任期内，扮演了很多角色——君主、外交官和文化象征。她不是每年一次出现在肖像画中，而是经常出现在公众视野里，正如女王自己所说：“我需要被人们看到，才能被他们相信。”

在英国国会开幕大典上，她会穿着隆重的服饰致辞。那件红色天鹅绒斗篷又长又重，需要几位侍从跟在身后抬起拖尾。每年6月，她都要在皇家赛马会上亮相。在为期5天的赛事中，城内谈资离不开女王的日间着装，有人甚至为“女王的帽子是什么颜色”而下注。到了夏季，白金汉宫会举办三场花园派对，荷里路德宫再加办一场，这也是英国公众与女王面对面的最好机会之一。

众多女王留下的影像不仅会被当作皇室风格来解读，也会被赋予某种政治意涵。1953年6月，伊丽莎白加冕仪式的全球转播仅英国就有2700万人收看，即便在小型黑白电视那种点状像素屏幕上，女王的礼服也能让人眼花缭乱。伊丽莎白要求设计师在白色缎面上以金银丝线描绘出英国和其他英联邦国家的象征，裙身上的装饰有爱尔兰的三叶草、加拿大的枫叶、新西兰的银蕨、澳大利亚的金合欢、南非的山龙眼等元素。毕竟，与这些国家保持友好关系是她的终身使命。

在位期间，女王对自己的政治倾向始终保持着沉默。作为一个极少能畅所欲言的人，她学会了通过挑选胸针等珠宝，来传递微妙的信息。解读女王珠宝的各种变体成了全球王室观察家的痴迷之事。2018年夏天，伊丽莎白在会见备受争议的美国前总统特朗普时，佩戴着奥巴马和米歇尔送给她的胸针，表面上看似保持中立，但这件象征友谊的物品，也透露出她对眼前人一种顽皮又狡黠的暗示。

如同历经15位英国首相和14位美国总统，伊丽莎白也见证过许多潮流，但她从未跟随过它们。曾有评论家指出，女王的两件套裙装是没有进化的时尚，着装风格缺乏现代感。但着装的一致性从另一种层面来说，是坚定和安心的象征。相比伊丽莎白一世的肖像将其描绘成穿着天鹅绒披着金丝斗篷的半神，维多利亚女王的形象如作家刘易斯·卡罗尔描述那般神似扑克红桃王后，让人想起“所有家庭教师的精华”，伊丽莎白二世是一位真正意义上的现代君主，在维系皇家习俗的同时，也小心翼翼地做出改变。更重要的是，在大多数君主政体都沦为仪式上的虚无，或已走向灭亡的时代，她重塑了君主权威，并将一种年长的、智慧的、慈爱又严格的母亲形象，转化为“国家”的拟人化表现。

\begin{quoteBox}{评价}
过去七十年风雨飘摇，充满了痛苦分裂和不可预知的动荡。尽管伊丽莎白被限制在她的身份里，禁锢在枯燥的宫廷仪式和习俗中，但她总能在最重要的时候，成为这个国家理想化的化身。

——英国历史学家西蒙·沙玛
\end{quoteBox}

\subsection*{传奇人生：王冠之下，一个女人的史诗}

当地时间2022年9月8日，英国女王伊丽莎白二世去世，享年96岁。作为当今在位时间最久、年纪最长的君主，她几乎经历了整个现当代史。无论早年战火纷飞时，还是晚年王室丑闻中，她总是语调平稳，波澜不惊。在头戴王冠的70年里，她带领英国王室朝着更加务虚的方向发展——远离政治纷争，把握微妙平衡，在若即若离中保持自身的影响力。

\textbf{有个性的淑女。}1926年4月21日凌晨，英国温莎家族迎来首位第三代成员。这名女婴被取名为“伊丽莎白”，很快便收获了整个英国王室的宠爱。彼时，谁也没将她与英国王位联系起来——她的父亲约克公爵是国王次子，而王位的第一继承人是她的伯父爱德华。

无须按王位继承人培养，伊丽莎白可以从容地接受“私人订制版”贵族淑女教育：家庭教师授课，课程除了写作、历史、法语，还包括钢琴、舞蹈等。在这种氛围中成长起来的伊丽莎白，也把“淑女范儿”刻进了骨髓。她14岁时，伦敦遭到纳粹德国的猛烈轰炸，即便是在炸弹来袭的夜晚，她也拒绝穿着睡衣去防空洞躲避，坚持换好衣服再撤。

当然，这位淑女并非全然不做“出格”的事。她会在伦敦水族馆的生日宴会上吹泡泡；在纵马奔驰时从不戴头盔。1942年，在“二战”尚未结束时，她还不顾父母反对，参加了英国军队的女子辅助服务团，担任汽车技工兼卡车司机。另外，出于天性和王室职责的需要，伊丽莎白在成为女王后也始终走在潮流的前沿。1974年，她成为第一个发送电子邮件的君主；2007年在YouTube上开通了专门的王室频道；2009年在Twitter上拥有了个人账号……

\textbf{守传统的女王。}英国前首相布莱尔曾评论称，伊丽莎白女王将传统视作维持君主制神秘和威严的关键。这些传统散布于各种王室礼仪与惯例中。1954年，已是女王的伊丽莎白在结束国事访问后登上返程的游艇，5岁的儿子查尔斯向她伸出小手，她却说道：“不，亲爱的，你得先等会儿。”说着，她先依次向船上的贵族和官员致意，然后才握了握儿子伸出的手。这件事给查尔斯留下了深刻的印象，他的传记作者安东尼·霍顿就将《不，亲爱的，你得先等会儿》作为书中一章的标题。

这种传统，是伊丽莎白的父亲教给她的。1936年12月11日，因执意迎娶曾离婚两次的辛普森太太，爱德华八世宣布逊位，约克公爵于茫然无措中成为新国王，即乔治六世。关于此事，乔治六世曾说：“我从来都不愿意发生这种事。我从没批阅过政务文件，我只是一个海军军官。”彼时的欧洲大陆，法西斯风潮迭起，英国国内人心浮动，随后“二战”爆发，战火波及英伦。不过，特殊事态下，乔治六世还是挑起了重担。他努力克服自己的口吃，频频发表演讲，到瓦砾堆中安抚和鼓励普通民众。而这一切，都被伊丽莎白看在眼里。

\subsection*{懂平衡的君主}

成为女王之后，伊丽莎白必须善用平衡之道——如果她过于矜持，就会失去对人民的亲和力；而如果过于平民化，又会丢掉女王这一身份的庄严。而与英国历代国王相比，伊丽莎白要履行的职责是多重的：君主、妻子和母亲，但可供参考的先例却几近于无。

伊丽莎白与菲利普亲王的婚恋故事曾被传为佳话，但这段感情也遭遇过冲击。伊丽莎白则在婚姻中展现了十足的隐忍和包容。而在1997年，戴安娜王妃因车祸去世于巴黎，当时流言纷纷，王室被架在火上。很少在公众面前表露私人情感的伊丽莎白，面对公关危机发表了一场电视讲话，说自己对儿媳的离去感到悲伤。丧礼当天，当戴安娜的棺木经过白金汉宫前，她低头鞠躬——身为英国的君主，女王是从不对任何人鞠躬的。最终，破天荒的举动消除了大半英国民众对王室的敌意，一个富有同情心的长辈形象通过电视传遍了世界。

伊丽莎白二世总是希望弥合分歧，凝聚家庭，用以助益国家。面对苏格兰对英联邦的离心力，她带着家人到苏格兰的巴尔莫勒尔庄园避暑，还让长子查尔斯和长孙威廉在苏格兰接受大学教育。在她的维系下，苏格兰人虽对英国政府十分不满，对王室却怀有敬意。

\begin{quoteBox}{名言}
我在你们所有人面前宣誓，此生无论短暂或漫长，我都将会为英联邦的全体人民、为我们同属的帝国家庭鞠躬尽瘁，奉献一生。

——1947年4月21日，伊丽莎白在自己21岁生日当天发表演讲
\end{quoteBox}

\begin{quoteBox}{名言}
对许多人来说，这是一个非常幸福和愉悦的时刻，但对那些失去亲人的人来说，圣诞节可能会很难熬的。特别是今年，我十分理解这一点。

——由于丈夫菲利普亲王于当年离世，伊丽莎白2021年圣诞致辞的主题是“向亲爱的菲利普致敬”。这也成为英国女王的最后一次圣诞致辞。
\end{quoteBox}

\subsection*{与昨日世界一起离开}

自丈夫菲利普亲王去世后，伊丽莎白的身体一天不如一天，一向不服老的她，开始在公开场合使用拐杖，还在医生的建议下戒了酒。然而，就在这个节骨眼上，儿子查尔斯还在给她“添堵”——其慈善机构因涉嫌“买卖爵位”遭到英国警方调查，引发舆论哗然。

查尔斯曾多次陷入舆论漩涡，此外，英国王室也经历过数次声誉危机。1992年，女王的几个孩子全部成为英国小报追逐的对象，王室被八卦化、丑角化。同年，王室行宫温莎城堡失火，损失惨重。在年底的演讲中，伊丽莎白自嘲说：“这一年的这些事让我不是很开心。”只是，大多数时候，她都选择将烦闷藏在心底。

不少人都认为，多年以来，英国君主制的维护主要在于女王一人，如果失去“女王保护”，英国王室内部的变数将层出不穷。然而，天下终究没有不散的筵席：96岁的人生，70年的王位，伊丽莎白从“二战”后开始，伴随英国进入了21世纪，见证了英国影响力的“江河日下”，看到了英国脱欧的“黑天鹅”，也兑现了自己年轻时的誓言。

\subsection*{另眼相看：伊丽莎白之后，英国王室何去何从？}
\auth{文/宋欣}

伊丽莎白女王离开了，英国有了新国王查尔斯三世，但人们还是会问：英国王室还能存在多久？

今年早些时候，伊丽莎白庆祝自己登基70周年，彼时她的人气十分高涨。据权威机构统计，英国有86\%的民众对“伊丽莎白作为女王的工作非常满意”。这个问题也问得很有意思，它并没有问“你是否对女王满意”，而是“对女王的工作是否满意”，这背后的含意是：虽然女王/国王是一个与生俱来的头衔，但是在当今它更多地被看成一份工作。既然是一份工作，那就意味着，你做得好可以留任，但做得不好，也可能会被“裁掉”。

英国王室的去留其实早已成为大家多年来讨论的焦点。1952年，当25岁的伊丽莎白继承王位的时候，她就应该明白，君主早已失去了实质性的政治权利，自己所要努力的就是让这个王位显得不那么苍白，维护王位以及王室的存在理由。

在伊丽莎白的“统治期间”，她无时无刻不注重将个人魅力和职业结合起来，确保自己平易近人的形象，但她又不能让外界走得太近以至于伤害到王室本身的利益。毕竟，王室是整个英国最富有的家族，伊丽莎白自己拥有超过6亿美元的净资产。而在经济危机、通货膨胀危机、能源危机的三重重压之下，英国民众生活水平急剧下降，不少英国普通民众甚至没有工资保障。

那么，王室能否长期存续？至少英国民众已经给出不容乐观的回答。根据2022年最新民调，有79\%的民众认为10年之后还会有王室，只有42\%的民众认为50年后还会有王室，只有不到三成的英国民众认为100年后王室还会存在。

这并不令人奇怪，从本质上来说，目前英国对于其他英联邦国家的影响力早已不复从前。澳大利亚、加拿大等英联邦国家随着自身在国际贸易领域的突飞猛进，国力不断上升，已经逐渐主导了和英国之间的外交关系。另一方面，英国王室如今的依存更多是靠伊丽莎白个人的魅力。但在信息产业高度发达的今天，在媒体无孔不入的当下，在新闻被不断娱乐化的全球，王室的“体面”不断丢失。从戴安娜事件开始，到后来的查尔斯婚外情，再到安德鲁王子丑闻、哈里和梅根事件，更多的人看到了王室存在逻辑的不合理。如今，随着伊丽莎白的离去，人们对于那个曾经辉煌一时的大不列颠的记忆也将逐渐模糊，留下的则是不断反省的英国社会。一如英国新任首相特拉斯26年前在牛津大学所倡议的：“没有人应当生下来就是统治者，人人生而平等。”

\begin{quoteBox}{Tips}
虽说英国实行君主立宪制，王室不掌握实权，但英国宪法的不成文属性以及复杂的政治惯例使得伊丽莎白二世发挥的实际作用要超过单纯的政坛吉祥物。纪录片《伊丽莎白二世——伟大的女王》记载，女王平均每年出席500个官方活动，每周例行与英国首相面谈，从未错过任何一次英联邦首脑会议。

而在英联邦以外的世界，伊丽莎白同样为和平做出了重要贡献。“二战”结束20周年，她出访联邦德国，标志着英德两国和解；1986年，她与丈夫在重阳节这天访问中国，并带来了400年前伊丽莎白一世写给万历皇帝的信作为国礼，成为历史上第一个访问中国的英国君主，释放出英国想与中国交好的强烈信号。
\end{quoteBox}

\subsection*{时评范本：伊丽莎白二世：作为女王和英国人的一生}
\auth{文/刘怡}

对于一个普通英国人来说，2022年的秋天无论从哪个角度看都是危机四伏的。猖獗的通货膨胀和生活必需品涨价潮已经席卷全国，30年来规模最大的罢工使部分地区的铁路陷入了间歇性“停摆”。而受乌克兰战争引发的能源市场动荡影响，经济学家普遍预测英国人入冬后的天然气消费账单将比12个月前上涨300\%。丽兹·特拉斯在9月6日被任命为过去6年里第四位英国首相，她从服饰风格到语言技巧都竭力模仿撒切尔夫人，但许多人怀疑她那项1800亿英镑的物价稳定计划能否按期兑现。

按照惯例，女王授权特拉斯组阁的流程会在伦敦的白金汉宫完成。不过这一次，47岁的新首相需要自行前往苏格兰阿伯丁郡的巴尔莫勒尔堡——自从今年2月从新冠肺炎感染中康复以来，伊丽莎白二世的健康状况一直不佳。而和过去几百年里一样，城堡中的生活依然是隐秘而封闭的，女王对这种神秘感的分寸也一向把握得很好。只是，她不会再像往常一样，在夏季结束后从苏格兰返回了。当地时间9月8日，英国王室正式发布了伊丽莎白二世于当天下午“在巴尔莫勒尔平静去世”的消息。

其实，英国王室在进入21世纪后的娴熟公关技巧有时会让人产生一种感觉，仿佛女王真的已经变为普通人，和公民一样拥有真实的喜怒哀乐、情绪波动，可以在儿孙绕膝中度过余年。这当然是一种错觉——女王从来都不是任何意义上的“普通人”。从出生开始，她没有踏进过真正的学校，童年女伴几乎都来自出生的家族，只能在一个很小的欧洲王室子弟圈子里选择自己的配偶。1952年继承英国君主之位时，她才25岁，喜欢骑马、狩猎和逗弄宠物狗。她和那些视她为统治者的人从来都不属于一个阶层，但她必须（至少在形式上）理解那一代人，否则王室本身都可能不复存在。

女王懂得利用大型仪式和公开出访的机会，为王室营造一种既有神秘感和等级特征，又不至于完全和公众脱节的形象。1957年，她开始把传统的圣诞节祝福改在电视直播中进行。而戴安娜事件之后，王室开始尝试主动利用公共舆论和社交媒体，消除大众与白金汉宫之间的隔阂。无论是威廉、哈里两位王子以及其他王室年轻成员成为娱乐新闻的主角，还是女王本人频频口出妙语，都成了主动利用曝光度的经典范例。

漫长的70年足以改变英国拥有的一切以及它所处的世界，也足以令英国君主制发生渐进、尽管并不是翻天覆地的变化。而这一切都是无法复制的。

\begin{quoteBox}{话题透视}
现今，英国王室在英国人心目中的崇高地位，很大程度上已不再来自旧日的“君权神授”思想，而是基于伊丽莎白二世个人道德所起到的榜样作用，是她在位70年所积累的声誉。莎士比亚曾在《亨利四世》中写下“欲戴王冠，必承其重”，伊丽莎白二世的人生确实承载了这样的重量。
\end{quoteBox}

\subsection*{立意指南}

\textbf{1. 勤政爱民，冷静隐忍。}英国女王令世人尊重的首要原因是她的“勤政爱民”。伊丽莎白在“二战”时便积极履行自己的职责：安慰民众、为士兵筹款，并加入了服务队。而在即位后，她频繁出访，利用当年英国王室的声望和殖民时代的共同记忆，不断维系英联邦的稳定，为英国王室乃至英国奉献了自己的一生。请看下面一则作文片段：

\begin{quoteBox}{作文片段}
作为现代社会中的一位“君主”，英国女王伊丽莎白二世恪守政治中立原则，不轻易表达自己的立场，但这并不代表她毫无作为。她致力维持联合王国的稳定，也经常外出访问；她宣扬平等主义，曾在访问乌干达时与该国总统跳舞，以此反对种族歧视。可以说，英国女王在英国乃至世界都有着很强的影响力。她鞠躬尽瘁，发挥着纽带的作用，为自己的国家奉献了一生，理应得到人民的爱戴。
\end{quoteBox}

\score{文段列举了英国女王所做的工作，用词准确且有分寸，不卑不亢，可谓十分得体。}

\textbf{2. 凝心聚力，指路明灯。}伊丽莎白二世在位时间长达70年，在漫长的履职生涯中，她发表了数千次演讲。其中很多颇具代表性的演讲都用智慧团结了英联邦，并成为民众的指路明灯，带领他们团结一致，走向未来。请看下面一则作文片段：

\begin{quoteBox}{作文片段}
英国女王伊丽莎白二世一直以自己的睿智和果决，如指路明灯般带领着联合王国走出困境。第二次世界大战期间，还是王储的伊丽莎白就展现出惊人的魄力，加入后勤部队为英军输送物资。英国“脱欧”闹剧愈演愈烈时，她下令阻止时任首相鲍里斯·约翰逊在未达成协议前退出欧盟，避免了不必要的混乱局面。而在新冠肺炎疫情肆虐英国之际，女王发表公开演说，呼吁人们共克时艰，渡过难关。伊丽莎白是一位优秀的领导者，以实际行动赢得了其子民的拥护。
\end{quoteBox}

\score{该文段最大的优点是材料翔实，让人信服，体现出深刻的道理——在困难的时候，越能站得高，就越能用自己的温暖来激励世人。}

\textbf{3. 王冠之下，亦是凡人。}伊丽莎白二世的优雅当然值得这个世界津津乐道，然而，作为女王，她始终不过是落后的君主制的代言人。与此同时，伊丽莎白竭力维系着君主制，也因家庭负面新闻的波及而遭到过少数社会舆论的非议。请看下面一则作文片段：

\begin{quoteBox}{作文片段}
伊丽莎白二世优雅随和，以自己的行动维护着英国王室的声誉和权威。在出席正式场合时，女王举止得体，充分体现着个人的修养。只是，伊丽莎白从未批评过英国的种族主义与殖民主义。为了英国或自身的利益，她也曾与独裁者来往，甚至用公共资金为自己深陷性侵丑闻的儿子提供法律辩护。她的君主之位是世袭的，在某种意义上代表着一种特权和无尽财富的傲慢。
\end{quoteBox}

\score{文段分两层，先写女王的优雅随和，用得体的行动和语言维护个人形象以及国家形象，接下来又辩证分析女王颇具争议的一面，令人深思。}

\subsection*{备考演练}
\auth{考场范本：垂拱而治，晓有为无为之道 · 广西桂林市国龙外国语学校高2107班 韦景严（指导老师：张华峰）}

\textbf{阅读下面的材料，按要求作文。}

英国女王伊丽莎白二世一生给世人留下的最深刻印象，是她强烈的责任感、对王位和臣民矢志不渝的献身精神。

有评论家说：“伊丽莎白二世是一位政治家，她最大的政治成就就是尽量不触碰政治。伊丽莎白二世担任女王的七十年间，是英国王室权力最小的时期，却也因此得到了英国民众最多的爱戴。”也有人评论说：“伊丽莎白二世的可贵之处，在于她从不参与具体政治，不让王室越俎代庖，不打破英国固有的政治平衡。但她也从不‘躺平’享乐，而是坚持着毕生工作的状态。”

以上材料引发了你怎样的思考？请写一篇文章，简要分析。要求：体裁自选，角度自定，题目自拟，不少于800字。

\textbf{考场范文}

英国女王伊丽莎白二世的去世引起了全世界的高度关注。其在位期间所积累的智慧，或许可用中国传统的道家思想来概括：“清静无为”“神人无功”等。在她作为女王的七十年间，英国顺利完成了数次转型，这离不开她“统而不治”式的“无为”态度。

“我无为，而民自化，我好静，而民自正。”女王的实践准则正好与之暗合。尽管英国的宪法中为君主保留了相当的权力，但是女王从不轻易使用，这也让她得到了“史上权力最小的女王”的称号。伊丽莎白即位的时代是一个创伤的时代，一个战后恢复并带着痛苦的牺牲与重燃的希望之时代。尽管如此，她仍然在变局中保持了高度的克制，对几乎所有首相的决策表示了信任，从不用个人观点干涉首相的政务。从撒切尔夫人到鲍里斯·约翰逊，尽管方针迥异，但女王都表现出极高的“适配性”。这也让英国王室在数次被卷入政治风波后最终都能平稳渡过，可谓达到了“垂拱而治”“帝力于我何有哉”的境界。

英国女王当然可以用自己的影响力“纵横捭阖”，但是深谙平衡之道的她也清楚自己的存在对英国保持温和传统的重要意义，王与相的平衡，君与民的平衡，都是她一生的重要课题，可说是“治大国如烹小鲜”。然而，“无为”不是无所作为，伊丽莎白在位期间充分地协调各方，苦心维持着国家的团结与统一。疫情期间，她充分发挥了国家象征的作用，号召本国民众共同抗疫，真正做到了“在其位，谋其政”的水准，也可谓“有所为有所不为”。

在位的七十年里，伊丽莎白经历了英国君主制度的一步步衰落，也见证了“大英帝国”一步步分解成松散的“英联邦”。但是，作为国家象征，她仍对英联邦的维持起到了重要作用。女王兢兢业业地出席每一个需要她的场合，充当着英国对外关系的润滑剂。她从不对他国内政指手画脚，而是通过个人魅力和优雅的礼节、风度，达成了英联邦国家的相对团结，维系着这些国家之间的平衡。

斯人已逝，但平衡有度的处世方式值得学习。做到当弟君子的中庸之道，达成“从心所欲，不逾矩”的张弛有度，也像英国女王一般知道“有所为有所不为”，方可宁静致远，以不变应万变。

\clearpage
\section*{10秒鲜素材}

\subsection*{1. 开启“上合时间”：奏合作之声，凝众行之力}

在喜迎丰收的仲秋之时，丝路古城、乌兹别克斯坦第二大城市撒马尔罕开启“上合时间”。9月14日至16日，中国国家主席习近平出席在撒马尔罕举行的上海合作组织成员国元首理事会第二十二次会议，并对哈萨克斯坦、乌兹别克斯坦两国进行国事访问。这是我国在党的二十大召开前夕开展的最重要的元首外交活动，充分体现了我国对上海合作组织和中哈、中乌关系的高度重视。

自2001年成立以来，上合组织不断发展壮大，已成为世界上幅员最广、人口最多，具有重要影响力和权威性的综合性区域组织，是完善国际秩序、推动地区和全球发展、维护成员国安全稳定不可或缺的建设性力量。而今年又适逢《上海合作组织宪章》签署20周年和《上海合作组织成员国长期睦邻友好合作条约》签署15周年。以撒马尔罕峰会为契机，践行“上海精神”，推动成员国共维稳定、共谋发展，才能为构建更加紧密的上海合作组织命运共同体，促进地区和平稳定、发展繁荣做出更大贡献。

上合组织已走过21个年头，成员国元首理事会第二十二次会议则为第三个十年的发展开了新局。在撒马尔罕峰会上，习近平主席同其他与会国家领导人一道，重点围绕上合组织框架内全方位合作、重大国际和地区问题等深入交换意见，凝聚合作新共识，谱写合作新篇章。面向未来，上合组织还将继续高举“上海精神”旗帜，走好团结合作之路、安危共担之路、开放融通之路、互学互鉴之路、公平正义之路，努力成为和谐、安宁、稳定、可持续发展的典范。

\begin{quoteBox}{立意标签}
1. 上合组织始终秉持互信、互利、平等、协商、尊重多样文明、谋求共同发展的“上海精神”，在国际形势深刻变化、新冠肺炎疫情延宕反复的背景下，为地区及世界和平发展注入了更多正能量。

2. 共促政治互信、共护安全稳定、共谋繁荣发展、共担国际道义，中国为上合组织的蓬勃发展、成员国互利合作贡献了不可或缺的力量。

3. 成员国之间政治、经济、文化、金融等领域的交流融通，是上合组织的重要内容。在这个平台上，各成员国都能找到彼此发展的“新支点”，共享“中国机遇”，实现共同繁荣。
\end{quoteBox}

\begin{quoteBox}{精彩片段}
和平与发展是时代主题，是不可抗拒的历史潮流。展望未来，只要沿着构建人类命运共同体的“人间正道”阔步前行，继续发出“上合声音”，贡献“上合方案”，我们坚信，上合组织各成员国必将绘就一幅各美其美、美美与共的时代新画卷。
\end{quoteBox}

\subsection*{2. “新版红绿灯”引热议：谣言何以引爆如此闹剧？}

近日，有媒体发布视频称“红绿灯新标准出台”，引发广泛关注。对于视频中提到的“取消红绿灯读秒”“将红黄绿三色灯改成九宫格组合”，不少网友表示不满。另据网友爆料，新版红绿灯设计者孙正良开直播仅10分钟便在网友抵制之下无奈关播。

不过，公安部交通管理科研所迅速辟谣指出，所谓“新国标红绿灯”并不新，已实施五年；目前各地在用道路交通信号灯符合现行标准的，不需要进行更换；复杂的“九宫格”式信号灯，是特殊组合的一种，“仅适用于极少数复杂路口”；而所谓的“新版红绿灯设计者直播”，系他人冒用孙正良名义。

复盘此事来龙去脉不难发现，事件起因是有媒体在对“新国标红绿灯”进行科普，受流量思维影响，自然不能免俗地要做“标题党”。进而，诸如“红绿灯取消读秒”“红绿灯变九宫格”等似是而非的说法，也便炮制出炉了。

一些网友会对“将红黄绿三色灯改成九宫格组合”的说法听之信之、焦虑不已，本质上传递出的，还是信心的缺乏。一是对公共治理水平缺乏信心，默认“化简为繁”的荒唐事真的会发生；二是对自己缺乏信心，担心看不懂那繁复的“九宫格”。我们需要坚信的是，包括交通管理在内，现代城市治理早有一套成熟的范式，“急转弯”“大切换”式的折腾，大概率是可以避免的。

\begin{quoteBox}{立意标签}
1. 新国标红绿灯总设计师孙正良曾荣获“有突出贡献中青年专家”称号，并两次获得国家科学技术进步奖。为了蹭热点而冒用孙正良名义进行直播，已侵害其合法权益。

2. 民生无小事，关乎千家万户的利益。越是便民、越是简单、越是直接，越能获得公众的理解和支持。

3. “新版红绿灯”最终被证实为谣言，这再次提醒我们，网络谣言必须整治，网络风气必须肃清。
\end{quoteBox}

\begin{quoteBox}{精彩片段}
一场“红绿灯”闹剧，再次照见了脆弱的网络生态。后真相时代，情绪大于事实，越来越多人沉溺其中，常常抓住个靶子就跟风宣泄情绪。在这样的背景下，新闻媒体、网络平台与相关部门责任重大，尤其应当好求证者、辟谣者，让真相“跑过”谣言。
\end{quoteBox}

\subsection*{3. 泸定抗震救灾：彰显不屈意志，谱写大爱乐章}

2022年9月5日，四川甘孜藏族自治州泸定县发生6.8级地震。在万众一心、众志成城的抗震救灾中，各方以责任与担当践行伟大的“抗震救灾精神”，汇聚动人至深的“中国力量”。

这力量是逢山开路、遇水架桥的担当。灾害发生后，来自解放军、森林消防、医护人员等多方的救援力量第一时间奔赴地震灾区，火速开展抢救伤员、灾后安置工作。在随时可能再次塌方的路上，他们冒着生命危险奋力抢通救援通道；在挡住去路的滚急流前，他们上演“飞夺泸定桥”式的“生命接力”。这力量是“一方有难，八方支援”的大爱。从中央组织部下拨2000万元“特殊党费”到地方通信、电力、铁路等部门开启应急模式，再到社会大众纷纷捐款捐物，来自全国各地各界的爱心如涓涓暖流汇入四川。

这力量也是自强不息、患难与共的责任。从幼儿园老师用臂膀护住孩子安全撤离到四川铁路部门第一时间扣停在途列车，确保客车和旅客安全，再到受灾群众积极参与现场救援，灾情出现后，震区各行各业答好了抗震减灾的考卷。

\begin{quoteBox}{写作参考}
今夏以来，从极端高温到水电供给困难，再到如今的地震，四川人民饱经考验，诸多不易，满是艰辛。但在磨难面前，他们展现出了惊人的韧性，而全国人民也守望相助，纷纷捐款捐物表示支持。在自然灾害面前，坚强善良的中华儿女团结一致，谱写了可歌可泣的感人诗篇。
\end{quoteBox}

\subsection*{4. 4部手机直播剪发：岂能用顾客权益换流量？}

2022年8月22日，武汉市民周女士在江汉区一家理发店剪发时，理发师未经允许使用4部手机直播了剪发过程，周女士要求停播却被怼“落伍”：“直播常态化了，又不是明星，普通老百姓怕啥？”随即，该新闻引发广泛关注。

基于门槛低、流量大、互动性强等特点，不少行业都选择通过直播提升关注度，但前提应该是征得顾客同意。而涉事理发师那句“又不是明星，普通老百姓怕啥”，恰恰反映出当下不少人对“普通老百姓”个人权利的忽视。近年来，直播花样百出，无辜“躺枪”者不在少数。此前，网红餐厅私自直播食客吃相、医生偷偷直播患者手术过程均引起公众批评，毕竟，未经允许私自直播不仅侵害了当事人的权益，还可能影响其正常生活。

直播盛行不是“越界”盛行的理由。在人人都有“摄像头”的时代，更应该了解摄像的禁区，这既是对自己的保护，也是对他人的尊重。

\begin{quoteBox}{写作参考}
说到底，理发师直播是为了营销、引流，但面对眼前的客人都不懂得尊重与珍惜，反而执意侵犯顾客隐私权和肖像权，又怎么能获得长远的发展呢？对服务业而言，真正能引流的，只有高超的技艺和一流的服务。而违法引流，只能是搬起石头砸自己的脚。
\end{quoteBox}

\subsection*{5. 我国首条跨海高铁全线贯通：绘就时代发展的壮美蓝图}

2022年8月30日，在厦门北站施工现场，随着最后一组500米长钢轨顺利铺设，我国首条跨海高铁——福厦高铁全线铺轨贯通。

福厦高铁设计时速350公里，是我国第一条真正意义上的海洋服役环境高铁，也是世界行车速度最快的跨海铁路。而震撼世界的“超级工程”背后，是“中国智造”这张金色名片不断焕新的灿烂光辉：国内自主研发的耐候钢新型材料，可“以锈制锈”长效防腐；抗拉强度高达2000兆帕的锌铝镀层钢丝，仅有7毫米的直径，却能承受80吨拉力……

福厦高铁通车后，福州、厦门两地也将实现“一小时生活圈”，厦门、泉州、漳州将形成“半小时交通圈”，为人民群众出行带来便利的同时，增强地区互联互通，拉动区域经济增长。

\begin{quoteBox}{写作参考}
从第一条高铁建成通车，到环沙漠铁路跑出“中国速度”，再到如今的跨海高铁全线贯通，每一项成就的背后，都是中国铁路人的智慧与汗水。面对重重困难，他们没有畏惧退缩，而是以“关关难过关关过”的斗志，相继攻克了道道技术难题。未来，向着星辰大海前进，中国铁路人必定能书写全新的中国故事。
\end{quoteBox}

\subsection*{6. 厚植网络文明土壤，共筑美好精神家园}

8月28日，2022年中国网络文明大会在天津国家会展中心正式开幕，大会发布《共建网络文明天津宣言》，提出“形成中国特色、中国风格、中国气派的网络文明新格局”，发出迈向数字文明新时代的中国新声。

近年来，互联网科技的飞速发展和网民数量的快速增加，给网络空间治理带来了新的挑战，各种低俗网络文化泛滥，网络诈骗令人防不胜防，网络谣言此起彼伏……为此，从2021年开始，国家网信办就集中开展了“清朗”系列专项行动，彰显出整治各类网络乱象的决心和信心。

党的十八大以来，党中央高度重视网络文明建设，把网络文明建设作为社会主义精神文明建设和网络强国建设的重要内容，政策举措有力有效，顶层设计日趋完善，推动网络空间正能量更加充盈、生态环境更加清朗。在全党全国喜迎党的二十大胜利召开之际，举办中国网络文明大会意义重大，通过展示我国网络文明建设成果，进一步增强全民自豪感、爱国情、凝聚力，必将在全国上下汇聚起共建网络文明的磅礴力量。

\begin{quoteBox}{写作参考}
网络世界不断繁荣的背后，离不开每个人主动成为网络文明的宣传员，凝聚合力形成崇德向善、见贤思齐的网络文明环境。如今，我们共绘网络文明美好画卷，共赴网络强国绚丽征程，必然能开启网络时代的全新篇章。
\end{quoteBox}

\subsection*{7. 取消5类全国性高考加分：为学子减负，促教育公平}

2022年9月15日，在教育部召开的“教育这十年”“1+1”系列第十二场新闻发布会上，相关负责人介绍，近年来，教育部持续清理规范高考加分，已明确取消了体育特长生、中学生奥林匹克竞赛、科技类竞赛、省级优秀学生、思想品德有突出事迹等5类全国性高考加分项目。

高考加分多年来饱受争议的重要原因在于，有加分项目就不可避免会出现人为操纵，不可避免会出现权力寻租，有钱有权的人有了比其他人更多的筹码，而没有“门路”的人在不知不觉中就被“弯道超车”。取消不必要的加分项目，是要让高考的招录不受到任何行政权力场外因素的影响，始终在法治的轨道上运行，始终在阳光下前行，最大力度地保证公开透明，充分发挥教育评价的积极导向作用，为所有考生公平地开启希望之门，为国家公正地选拔出优秀人才。

\begin{quoteBox}{写作参考}
过去一些年里，校外培训产业如火如荼，与此同时，学生的负担越来越重，家长的教育焦虑也不断加深。但学生上“兴趣班”未必是出于“兴趣”，而很可能是为了在升学时增加砝码。因此，无论“双减”还是教育部门清理规范高考加分，都既符合教育规律，也适应时代变化，能有效地给教育领域喧嚣浮躁的情绪降温，让教育回归教书育人的本质。
\end{quoteBox}

\subsection*{8. “网红路牌”违规被拆，“浪漫风”不能随性吹}

“我喜欢乐山，更喜欢这座城市的你”“我在宁波很想你”“想你的风还是吹到了重庆”……日前，一系列“浪漫风”路牌在多地出现，并掀起拍照打卡热。这种外形酷似道路标志的“网红路牌”其实多是附近商家私自设置的“引流”利器，已有不少因涉嫌违规被拆除。

道路标志是国家法定标志物，设置路标更是一件严肃的事情。正规路牌要依据国家标准统一设计制作，既具有正确指示道路的功能，又符合安全要求。由此可见，商家为了“引流”，将自行设置的标志牌立于道路两旁，把生意算盘打到地名标志上，不仅越界，还涉嫌违规。试想，如果街道上“网红路牌”林立，很可能对不清楚情况的路人、司机产生误导，更别提打卡聚集可能引发交通事故这样的惨剧。

商家在追求经济利益的同时，还要承担起相应的社会责任。如果营销行为对维持公共秩序不能起到正向作用，恐怕也留不住消费者的脚步。

\begin{quoteBox}{写作参考}
竖立在交通道路上的路牌承担着指路作用，不可随意设置。以假乱真的“网红路牌”越来越多，势必误导公众，更会影响正常的公共交通秩序。只是，“网红路牌”能成为打卡地，证明其某种程度上契合了一部分人的情感需求，也不必“一刀切”拆除。如何在保障城市正常运行的基础上正确处理“网红路牌”，考验着城市管理的智慧。
\end{quoteBox}

\subsection*{9. 卖5斤芹菜被罚6.6万元：执法需要力度，也要温度}

近日，国务院督查组在国务院“互联网+督查”平台上接到群众反映，陕西省榆林市一家个体户卖出5斤芹菜后，被市场监管部门以品质不合格为由罚款6.6万元。随即，督查组对此展开调查走访。而经过调查，当地市场监管部门承认“确实过罚不当”。

事实上，此案并非个例。督查组查阅榆林市市场监管局2021年以来食品类行政处罚台账发现，在针对小微市场主体的五十多起处罚中，罚款超过5万元的就有21起，而他们的案值只有几十或几百元。《中华人民共和国行政处罚法》规定，设定和实施行政处罚必须以事实为依据，与违法行为的事实、性质、情节以及社会危害程度相当。即便对食品安全问题“零容忍”，执法者也不能无视法律规定，滥收滥罚，要有严格执法的“冷面孔”，也应有服务社会的“热心肠”。

\begin{quoteBox}{写作参考}
“卖5斤芹菜被罚6.6万元”当事人进货7斤芹菜，2斤送检，最终卖出20元，扣除成本后盈利只有几元钱。这样的收入与巨额处罚形成强烈反差，于是引发公众同情。食品安全怎么重视都不为过，但重视不能等于重罚。如何让老百姓吃得放心，又做到监管合情合理合法，需要相关部门深入思考。
\end{quoteBox}

\subsection*{10. 《我的县长父亲》遭舆论“误伤”，我们应该反思什么？}

2022年9月18日，山东省德州市作家协会公布“廉洁文化主题文学作品征文获奖名单”，一篇题为《我的县长父亲》的散文获得一等奖。很快，该文就因题目与喜剧电影《夏洛特烦恼》中的获奖作文《我的区长父亲》相似，受到部分网友嘲讽，征文评选活动的公正性也因此受到质疑。争议发生后，德州作协一度删除了网上的获奖信息与文章链接，但舆论很快又发生了反转。

《我的县长父亲》的作者于忠东是一位退休老人，在文章中，她讲述了父亲于志明的一些感人事迹，一个甘于清贫、有情有义、无私奉献、造福一方的基层干部形象跃然纸上。看完全文的网友纷纷表示，这篇作品获得一等奖实至名归。作为评选单位，德州作协却迫于压力删除了获奖信息，一删了之看起来简单，但不仅没有平息争议，反而给人“此地无银三百两”的错觉。

面对复杂多变的舆情，任何选拔类活动的组织者都不能被一时风向裹挟带偏，只要标准、程序正当，就经得起检验和审视。而网友们也需保持理性，不被先入为主的情绪或偏见所误导。

\begin{quoteBox}{写作参考}
假如《我的县长父亲》这篇文章更换为其他标题，可能并不会引发关注，同理，如果网友读过文章，恐怕也不会盲目质疑。种种逻辑共同折射出当下网络氛围的浮躁，这也提醒我们，在互联网时代更要保持理性，发言前不妨先求证，以此共同维护清朗有序的网络环境。
\end{quoteBox}
